\section{Experiments}
\label{sec:experiments}
\begin{table*}[t]
\centering
\caption{University-1652 retrieval performance (\%) and model efficiency. FPS is measured on Jetson Orin NX.}
\label{tab:university}
\footnotesize
\setlength{\tabcolsep}{3pt}
\begin{tabular*}{\textwidth}{@{\extracolsep{\fill}}lcrrrrrrrr@{}}
\toprule
\raisebox{-4.5pt}[0pt][0pt]{Method} & \raisebox{-4.5pt}[0pt][0pt]{Venue / Year} & \raisebox{-4.5pt}[0pt][0pt]{Input} & \raisebox{-4.5pt}[0pt][0pt]{Params (M)$\downarrow$} & \raisebox{-4.5pt}[0pt][0pt]{FLOPs (G)$\downarrow$} & \raisebox{-4.5pt}[0pt][0pt]{FPS$\uparrow$} & \multicolumn{2}{c}{D$\rightarrow$S} & \multicolumn{2}{c}{S$\rightarrow$D} \\
 & & & & & & R@1$\uparrow$ & AP$\uparrow$ & R@1$\uparrow$ & AP$\uparrow$ \\
\midrule
LPN~\cite{wang2022lpn}$^{r}$ & TCSVT 2022 & 256 & 27.71 & 16.28 & 33.28 & 75.77 & 79.09 & 88.02 & 76.01 \\
MuSe-Net~\cite{wang2024musenet}$^{r}$ & PR 2024 & 256 & 26.60 & 25.88 & 19.98 & 74.73 & 78.09 & 87.87 & 75.52 \\
Sample4Geo~\cite{deuser2023sample4geo}$^{r}$ & ICCV 2023 & 384 & 87.57 & 90.83 & 24.06 & 92.66 & 93.81 & 95.15 & 91.39 \\
MEAN~\cite{chen2025mean} & TGRS 2025 & 384 & 27.82 & 26.46 & 40.74 & 93.55 & 94.53 & 96.01 & 92.08 \\
CAMP~\cite{wu2024camp}$^{r}$ & TGRS 2024 & 384 & 87.71 & 90.83 & 18.73 & 94.42 & 95.35 & 96.15 & 92.72 \\
DAC~\cite{xia2024dac}$^{r}$ & TCSVT 2024 & 384 & 87.57 & 90.83 & 18.76 & \textbf{94.68} & \textbf{95.51} & \textbf{96.43} & \textbf{93.80} \\
CCR~\cite{du2024ccr}$^{r}$ & TCSVT 2024 & 384 & 121.68 & 90.96 & 10.13 & 92.95 & 94.16 & 95.72 & 92.55 \\
MCCG~\cite{shen2024mccg}$^{r}$ & TCSVT 2024 & 256 & 29.00 & 11.77 & 60.05 & 89.59 & 91.33 & 94.44 & 89.62 \\
SpikeViMFormer-S~\cite{chen2026spikevimformer}$^{r}$ & TCSVT 2026 & 384 & 18.63 & 50.12 & 32.49 & 88.00 & 89.96 & 92.30 & 87.11 \\
InfoGeo-RD~\cite{zhang2026infogeo}$^{r}$ & ICML 2026 & 448 & 91.57 & 223.50 & 18.24 & 89.82 & 91.43 & 95.29 & 90.32 \\
DepthGeoLoc~\cite{sun2026depthgeoloc} & TGRS 2026 & 252 & 28.57 & 11.63 & -- & 92.75 & 93.82 & 94.58 & 91.56 \\
MobileGeo~\cite{sun2025mobilegeo} & arXiv 2025 & 224 & 28.57 & \textbf{4.45} & -- & 93.87 & 94.83 & 95.72 & 92.57 \\
\midrule
\textbf{EfficientGeo} & Ours & 256 & 13.62 & 6.07 & 137.67 & 93.20 & 94.28 & 95.15 & 92.34 \\
\textbf{EfficientGeo + Rep.} & Ours & 256 & \textbf{13.54} & 6.01 & \textbf{214.34} & 93.14 & 94.23 & 95.29 & 92.32 \\
\bottomrule
\end{tabular*}
\par\smallskip{\footnotesize\raggedright $^{r}$: reproduced results. Other results are paper-reported. DepthGeoLoc: student. MobileGeo: without MSRM. Rep.: reparameterized. --: unavailable.\par}
\end{table*}



\subsection{Experimental Setup}
\label{sec:experimental_setup}

\paragraph{Datasets.}
To assess the cross-dataset generalization of EfficientGeo,
we train the model on University-1652~\cite{zheng2020university}
and evaluate it on SUES-200~\cite{zhu2023sues} and
GTA-UAV~\cite{ji2025game4loc} in a zero-shot setting. On
SUES-200, we report results separately at four flight
heights of 150, 200, 250, and 300\,m. On GTA-UAV, we
report results on both the same-area and cross-area test
splits. Further details of the three datasets are provided
in Section B of the supplementary material.

\begin{table*}[t]
\centering
\caption{Zero-shot transfer from University-1652 to SUES-200. Each cell reports R@1/AP (\%) at the indicated height.}
\label{tab:transfer}
\scriptsize
\setlength{\tabcolsep}{1.5pt}
\begin{tabular*}{\textwidth}{@{\extracolsep{\fill}}lc*{8}{c}@{}}
\toprule
& & \multicolumn{4}{c}{D$\rightarrow$S} & \multicolumn{4}{c}{S$\rightarrow$D} \\
\cmidrule(lr){3-6}\cmidrule(lr){7-10}
Method & Params (M)$\downarrow$ & 150 m & 200 m & 250 m & 300 m & 150 m & 200 m & 250 m & 300 m \\
\midrule
LPN~\cite{wang2022lpn}$^{r}$ & 27.71 & 41.28/49.24 & 53.00/59.83 & 59.58/65.59 & 63.98/69.61 & 36.25/34.20 & 46.25/45.69 & 57.50/53.17 & 65.00/59.02 \\
MuSe-Net~\cite{wang2024musenet}$^{r}$ & 26.60 & 65.63/70.04 & 77.13/80.02 & 81.38/83.98 & 83.65/85.98 & 75.00/63.78 & 80.00/73.78 & 85.00/79.45 & 86.25/81.06 \\
Sample4Geo~\cite{deuser2023sample4geo}$^{r}$ & 87.57 & 74.18/78.05 & 83.93/86.75 & 88.38/90.51 & 91.90/93.43 & 85.00/78.90 & 90.00/87.81 & 95.00/92.62 & 92.50/92.88 \\
MEAN~\cite{chen2025mean} & 27.82 & 81.73/87.72 & 89.05/91.00 & 92.13/93.60 & 94.63/95.76 & 91.25/81.50 & 96.25/89.55 & 95.00/92.36 & 96.25/94.32 \\
CAMP~\cite{wu2024camp}$^{r}$ & 87.71 & 78.90/82.37 & 86.83/89.27 & 91.95/93.63 & 95.68/96.65 & 87.50/78.98 & 95.00/87.04 & 95.00/91.05 & 96.25/93.44 \\
DAC~\cite{xia2024dac}$^{r}$ & 87.57 & 76.65/80.57 & 86.48/89.01 & 92.98/94.20 & 94.53/95.45 & 87.50/79.87 & 96.25/88.98 & 95.00/92.81 & 96.25/93.99 \\
CCR~\cite{du2024ccr}$^{r}$ & 121.68 & 73.33/77.72 & 86.73/89.05 & 93.23/94.43 & 95.53/96.32 & 86.25/75.31 & 92.50/88.36 & 95.00/93.93 & \textbf{97.50}/95.80 \\
MCCG~\cite{shen2024mccg}$^{r}$ & 29.00 & 78.53/82.29 & 87.63/89.97 & 92.65/94.13 & 94.95/96.06 & 82.50/75.22 & 88.75/84.87 & 93.75/90.23 & 95.00/92.36 \\
SpikeViMFormer-S~\cite{chen2026spikevimformer}$^{r}$ & 18.63 & 76.05/79.95 & 83.88/86.67 & 87.80/89.90 & 89.25/91.13 & 85.00/74.63 & 88.75/82.90 & 90.00/86.16 & 88.75/87.97 \\
InfoGeo-RD~\cite{zhang2026infogeo}$^{r}$ & 91.57 & \textbf{91.88}/\textbf{93.15} & \textbf{95.40}/\textbf{96.11} & \textbf{96.83}/\textbf{97.26} & 96.68/97.18 & \textbf{97.50}/\textbf{89.17} & 96.25/\textbf{92.99} & \textbf{98.75}/\textbf{95.41} & 96.25/\textbf{96.04} \\
MobileGeo~\cite{sun2025mobilegeo} & 28.57 & -- & -- & -- & -- & 92.50/83.81 & \textbf{97.50}/91.75 & \textbf{98.75}/94.59 & \textbf{97.50}/\textbf{96.04} \\
\midrule\textbf{EfficientGeo} & 13.62 & 86.13/88.73 & 94.10/95.22 & 96.53/97.19 & 97.58/98.05 & 92.50/83.88 & \textbf{97.50}/91.56 & \textbf{98.75}/94.16 & 96.25/95.00 \\
\textbf{EfficientGeo + Rep.} & \textbf{13.54} & 86.15/88.73 & 94.03/95.16 & 96.55/97.20 & \textbf{97.60}/\textbf{98.07} & 92.50/83.65 & \textbf{97.50}/91.58 & \textbf{98.75}/94.17 & \textbf{97.50}/95.04 \\
\bottomrule
\end{tabular*}
\par\smallskip{\footnotesize\raggedright $^{r}$: reproduced results. Other results are paper-reported. --: unavailable.\par}
\end{table*}

\begin{table}[!b]
\centering
\caption{Zero-shot drone-to-satellite retrieval on GTA-UAV. R@1 and AP are in \%. Dis@1 is in meters.}
\label{tab:gta_transfer}
\scriptsize
\setlength{\tabcolsep}{1.5pt}
\resizebox{\columnwidth}{!}{%
\begin{tabular}{lccccc}
\toprule
\multicolumn{6}{c}{Same-area} \\
Method & Params (M) & R@1$\uparrow$ & AP$\uparrow$ & SDM@3$\uparrow$ & Dis@1$\downarrow$ \\
\midrule
LPN~\cite{wang2022lpn}$^{r}$ & 27.71 & 12.78 & 18.96 & 0.2968 & 2258.60 \\
MuSe-Net~\cite{wang2024musenet}$^{r}$ & 26.60 & 25.88 & 33.28 & 0.4147 & 1634.71 \\
Sample4Geo~\cite{deuser2023sample4geo}$^{r}$ & 87.57 & 29.28 & 38.18 & 0.4770 & 1363.48 \\
MEAN~\cite{chen2025mean}$^{r}$ & 27.82 & 30.87 & 38.83 & 0.4730 & 1486.16 \\
CAMP~\cite{wu2024camp}$^{r}$ & 87.71 & 34.71 & 43.62 & 0.5071 & 1326.12 \\
DAC~\cite{xia2024dac}$^{r}$ & 87.57 & 34.24 & 43.45 & 0.5005 & 1374.89 \\
CCR~\cite{du2024ccr}$^{r}$ & 121.68 & 32.36 & 40.78 & 0.4813 & 1308.37 \\
MCCG~\cite{shen2024mccg}$^{r}$ & 29.00 & 35.49 & 43.85 & 0.4742 & 1414.49 \\
SpikeViMFormer-S~\cite{chen2026spikevimformer}$^{r}$ & 18.63 & 23.06 & 30.54 & 0.4089 & 1704.76 \\
InfoGeo-RD~\cite{zhang2026infogeo}$^{r}$ & 91.57 & 32.21 & 40.53 & 0.5112 & \textbf{1224.02} \\
\midrule
\textbf{EfficientGeo} & 13.62 & \textbf{40.37} & \textbf{48.16} & \textbf{0.5139} & 1282.56 \\
\textbf{EfficientGeo + Rep.} & \textbf{13.54} & 40.17 & 47.99 & 0.5123 & 1288.51 \\
\bottomrule
\end{tabular}}\par\smallskip
\resizebox{\columnwidth}{!}{%
\begin{tabular}{lccccc}
\toprule
\multicolumn{6}{c}{Cross-area} \\
Method & Params (M) & R@1$\uparrow$ & AP$\uparrow$ & SDM@3$\uparrow$ & Dis@1$\downarrow$ \\
\midrule
LPN~\cite{wang2022lpn}$^{r}$ & 27.71 & 12.83 & 20.04 & 0.3230 & 2050.10 \\
MuSe-Net~\cite{wang2024musenet}$^{r}$ & 26.60 & 25.82 & 34.10 & 0.4382 & 1389.22 \\
Sample4Geo~\cite{deuser2023sample4geo}$^{r}$ & 87.57 & 31.11 & 40.33 & 0.5115 & 1051.62 \\
MEAN~\cite{chen2025mean}$^{r}$ & 27.82 & 29.99 & 38.65 & 0.5005 & 1185.61 \\
CAMP~\cite{wu2024camp}$^{r}$ & 87.71 & 36.09 & 45.43 & 0.5483 & 928.02 \\
DAC~\cite{xia2024dac}$^{r}$ & 87.57 & 36.23 & 45.73 & 0.5335 & 1052.64 \\
CCR~\cite{du2024ccr}$^{r}$ & 121.68 & 32.46 & 41.47 & 0.5100 & 1101.47 \\
MCCG~\cite{shen2024mccg}$^{r}$ & 29.00 & 36.89 & 46.21 & 0.5068 & 1109.27 \\
SpikeViMFormer-S~\cite{chen2026spikevimformer}$^{r}$ & 18.63 & 20.99 & 29.24 & 0.4381 & 1364.73 \\
InfoGeo-RD~\cite{zhang2026infogeo}$^{r}$ & 91.57 & 33.23 & 42.05 & 0.5345 & 962.25 \\
\midrule
\textbf{EfficientGeo} & 13.62 & 43.48 & 51.99 & \textbf{0.5681} & \textbf{892.75} \\
\textbf{EfficientGeo + Rep.} & \textbf{13.54} & \textbf{43.60} & \textbf{52.03} & 0.5673 & 894.28 \\
\bottomrule
\end{tabular}}
\par\smallskip{\footnotesize\raggedright $^{r}$: our reproduced retrieval results.\par}
\end{table}

\begin{table}[t]
\centering
\caption{Query-encoding efficiency on Jetson Orin NX. Rep.: RepViT structural reparameterization.}
\label{tab:jetson}
\footnotesize
\setlength{\tabcolsep}{2pt}
\resizebox{\columnwidth}{!}{%
\begin{tabular}{lcccc}
\toprule
Method & \shortstack{Latency\\(ms)$\downarrow$} & \shortstack{Peak alloc.\\(MiB)$\downarrow$} & \shortstack{Power\\(W)$\downarrow$} & \shortstack{Energy\\(J/image)$\downarrow$} \\
\midrule
LPN~\cite{wang2022lpn} & 30.04$\pm$0.01 & 241.32 & 24.72 & 0.7454 \\
MuSe-Net~\cite{wang2024musenet} & 50.05$\pm$0.07 & 331.36 & 25.10 & 1.2648 \\
Sample4Geo~\cite{deuser2023sample4geo} & 41.57$\pm$0.16 & 362.11 & 23.63 & 0.9789 \\
MEAN~\cite{chen2025mean} & 24.55$\pm$0.02 & 169.86 & 21.88 & 0.5346 \\
CAMP~\cite{wu2024camp} & 53.39$\pm$0.18 & 376.43 & 23.33 & 1.2412 \\
DAC~\cite{xia2024dac} & 53.29$\pm$0.06 & 396.55 & 22.97 & 1.2205 \\
CCR~\cite{du2024ccr} & 98.69$\pm$0.06 & 632.56 & 27.11 & 2.6693 \\
MCCG~\cite{shen2024mccg} & 16.65$\pm$0.01 & 148.11 & 25.65 & 0.4241 \\
SpikeViMFormer-S~\cite{chen2026spikevimformer} & 30.78$\pm$0.04 & 100.55 & 21.11 & 0.6507 \\
InfoGeo-RD~\cite{zhang2026infogeo} & 54.82$\pm$0.03 & 607.77 & 24.78 & 1.3557 \\
\midrule
DINOv3 ViT-S/16 & 6.65$\pm$0.02 & 69.02 & 20.46 & 0.1345 \\
DINOv3 ConvNeXt-T & 9.85$\pm$0.01 & 79.85 & 20.04 & 0.1960 \\
\midrule
\textbf{EfficientGeo} & 7.264$\pm$0.021 & 27.82 & \textbf{17.29} & 0.1242 \\
\textbf{EfficientGeo + Rep.} & \textbf{4.665}$\pm$0.014 & \textbf{27.35} & 20.34 & \textbf{0.0932} \\
\bottomrule
\end{tabular}%
}
\end{table}
\paragraph{Implementation details.}
University-1652 provides training and test splits but no
separate validation split~\cite{zheng2020university}.
As in the released implementations of
Sample4Geo~\cite{deuser2023sample4geo},
DAC~\cite{xia2024dac}, and
MEAN~\cite{chen2025mean}, we evaluate on its test split
during training and save checkpoints according to test
performance. The selected student is evaluated on SUES-200
and GTA-UAV in a zero-shot setting. The main retrieval tables
report results for both the unreparameterized and reparameterized student.Further training and optimization details are provided in Section A.2 of the supplementary material.

On University-1652 and SUES-200, we report Recall@1
(R@1) and average precision (AP) for both
drone-to-satellite (D$\rightarrow$S) and
satellite-to-drone (S$\rightarrow$D) retrieval.
Following the evaluation implementations of
Sample4Geo~\cite{deuser2023sample4geo},
DAC~\cite{xia2024dac}, and
MEAN~\cite{chen2025mean}, we exclude gallery images
whose identities are absent from the query set when
computing these metrics.

For GTA-UAV, we follow the official evaluation
protocol~\cite{ji2025game4loc} and report R@1, AP,
SDM@3, and Dis@1 for D$\rightarrow$S retrieval.
\subsection{Cross-View Retrieval Performance}
\label{sec:retrieval_performance}



\paragraph{University-1652.}
EfficientGeo achieves 93.20\% D$\rightarrow$S and 95.15\% S$\rightarrow$D R@1 with 13.62\,M parameters (Table~\ref{tab:university}). The lightweight SpikeViMFormer-S achieves 88.00\% and 92.30\% in the two directions with 18.63\,M parameters. DAC achieves 94.68\% and 96.43\% with 87.57\,M parameters. These results show that the compact student retains strong source-domain discrimination with substantially fewer parameters.

\paragraph{Zero-shot transfer to SUES-200.}
EfficientGeo achieves an average D$\rightarrow$S R@1 of 93.58\% across the four flight heights, outperforming all compared methods except InfoGeo-RD (Table~\ref{tab:transfer}). InfoGeo-RD achieves approximately 95.20\% with 91.57\,M parameters, compared with 13.62\,M for EfficientGeo. At 300\,m, EfficientGeo reaches 97.58\% R@1, compared with 96.68\% for InfoGeo-RD. The lightweight SpikeViMFormer-S achieves approximately 84.25\% average R@1. These results indicate that the compact student retains transferable representations across changes in flight height.

\paragraph{Zero-shot transfer to GTA-UAV.}
EfficientGeo achieves 40.37\%/48.16\% and 43.48\%/51.99\% R@1/AP on the same-area and cross-area splits, respectively (Table~\ref{tab:gta_transfer}). In comparison, MCCG achieves 35.49\%/36.89\% R@1 on the two splits, and the lightweight SpikeViMFormer-S achieves 23.06\%/20.99\% with 18.63\,M parameters. EfficientGeo also exceeds InfoGeo-RD in R@1 while using substantially fewer parameters (13.62\,M vs.\ 91.57\,M). These results suggest that the distilled student retains strong cross-region generalization under a compact model budget. On GTA cross-area, EfficientGeo achieves an SDM@3 of 0.5681 and a Dis@1 of 892.75\,m, compared with 0.5068 and 1109.27\,m for MCCG. On the cross-area split, InfoGeo-RD and Sample4Geo obtain Dis@1 values of 962.25\,m and 1051.62\,m, respectively, both higher than EfficientGeo's. On same-area retrieval, InfoGeo-RD achieves a lower Dis@1 of 1224.02\,m, compared with 1282.56\,m for EfficientGeo.


\subsection{Deployment Efficiency on Jetson Orin NX}
\label{sec:deployment_efficiency}
\paragraph{Measurement setup.}
We measure query encoding on Jetson Orin NX 16\,GB in MAXN\_SUPER mode with fixed clocks, batch size one, CUDA Graph, and each model's native resolution and precision. After warm-up, each of three rounds measures 1,024 queries. Timing includes copying the GPU-resident input into the CUDA Graph buffer, encoding, normalization, and synchronization. We report the average over three rounds. Further details of the deployment measurements are provided in Section C of the supplementary material.

\paragraph{Deployment performance.}On Jetson Orin NX, reparameterized EfficientGeo encodes a query in 4.665\,ms, with a peak GPU allocation of 27.35\,MiB and an energy cost of 0.0932\,J per query (Table~\ref{tab:jetson}). Under the same input resolution and BF16 precision, these values are 29.8\%, 60.4\%, and 30.7\% lower, respectively, than those of the fine-tuned DINOv3 ViT-S/16. RepViT's native structural reparameterization~\cite{wang2024repvit} reduces latency by 35.8\% and energy per query by 25.0\% relative to the unreparameterized student, while power increases from 17.29\,W to 20.34\,W. The lightweight SpikeViMFormer-S has a query latency of 30.78\,ms and a peak GPU allocation of 100.55\,MiB on the same device. Reparameterized EfficientGeo also surpasses it in zero-shot R@1 on SUES-200 and both GTA-UAV splits. Together, these results show that EfficientGeo combines strong cross-region generalization with efficient onboard inference. Per-model inference settings and structural reparameterization results are provided in Sections C and D of the supplementary material, respectively.

\subsection{Ablation Studies}
\label{sec:ablations}
\begin{table}[!b]
\centering
\caption{Stage-I ablations of intermediate-teacher training. Cells report D$\rightarrow$S R@1/AP (\%). SUES-200 averages four heights. Cal.: feature calibration. U1652: University-1652. GTA-S/C: GTA-UAV same-/cross-area.}
\label{tab:ablation_middle}
\setlength{\tabcolsep}{1.5pt}
\fontsize{7.5}{9}\selectfont
\begin{tabular*}{\linewidth}{@{\extracolsep{\fill}}lcccc@{}}
\toprule
Variant & U1652 & SUES-200 & GTA-S & GTA-C \\
\midrule
W/o distillation & 94.68/95.61 & 92.91/94.26 & 41.97/51.15 & 45.85/55.12 \\
\midrule
\multicolumn{5}{@{}l}{\textit{Distillation targets}} \\
Similarity KL & 94.81/95.70 & 92.80/94.13 & 41.50/50.78 & 45.22/54.68 \\
\addlinespace[2pt]
HND ($M$-selected) & 95.26/96.08 & 93.18/94.42 & 40.87/49.78 & 44.52/53.84 \\
\addlinespace[2pt]
TG-HND & 95.62/96.42 & 94.63/95.63 & 43.16/52.32 & 47.27/56.55 \\
\midrule
\multicolumn{5}{@{}l}{\textit{Alternative objective (with DG-SAM)}} \\
ADSD & 95.26/96.08 & 93.11/94.36 & 40.08/48.90 & 43.07/52.33 \\
\midrule
\multicolumn{5}{@{}l}{\textit{TG-HND extensions}} \\
+ Cal. & 95.80/96.55 & 94.39/95.40 & 42.70/51.78 & 46.48/56.12 \\
\addlinespace[2pt]
+ DG-SAM & 95.93/96.65 & \textbf{95.53}/\textbf{96.32} & 44.58/53.38 & 47.60/57.18 \\
\addlinespace[2pt]
+ Cal. + DG-SAM & \textbf{96.15}/\textbf{96.84} & 95.49/96.26 & \textbf{46.04}/\textbf{54.70} & \textbf{49.65}/\textbf{58.71} \\
\bottomrule
\end{tabular*}

\end{table}
\paragraph{Stage I: teacher-to-intermediate distillation.}
Table~\ref{tab:ablation_middle} starts from an intermediate teacher trained without distillation. On GTA cross-area, similarity KL reduces D$\rightarrow$S R@1 from 45.85\% to 45.22\%, whereas TG-HND raises it to 47.27\%. On SUES-200 and GTA same-area, TG-HND reaches 94.63\% and 43.16\% R@1, compared with 92.80\% and 41.50\% for KL. When the foundation teacher still supplies the margin targets but the intermediate teacher selects hard negatives, GTA same-area and cross-area R@1 fall from 43.16\%/47.27\% to 40.87\%/44.52\%. This controlled comparison highlights the value of foundation-guided negative selection for transferring discriminative rankings across the capacity gap. With DG-SAM fixed, TG-HND reaches 44.58\%/47.60\% R@1 on the two GTA splits, compared with 40.08\%/43.07\% for ADSD. This supports teacher-guided ranking transfer for the tested same-architecture teacher pair.

Feature calibration and DG-SAM further strengthen the intermediate teacher when used together. Relative to TG-HND alone, their combination raises GTA same-area and cross-area R@1 from 43.16\%/47.27\% to 46.04\%/49.65\%. With DG-SAM fixed, calibration raises GTA cross-area R@1 from 47.60\% to 49.65\%, while adding DG-SAM to the calibrated model increases it from 46.48\% to 49.65\%. Calibration alone lowers this score from 47.27\% to 46.48\%, indicating that its benefit depends on the optimization setting. Detailed ablation results are provided in Section E of the supplementary material.

\paragraph{Stage II: cross-architecture distillation.}
Table~\ref{tab:ablation_student} examines how subspace supervision and DG-SAM affect transfer to the RepViT student. Compared with training without distillation, teacher-subspace supervision improves zero-shot retrieval on GTA-UAV, indicating that the intermediate teacher provides transferable guidance. The benefit of the full ADSD objective depends on optimization. Without DG-SAM, ADSD does not outperform TS-only on GTA cross-area, whereas with DG-SAM it raises R@1 from 40.21\% to 43.48\%. This comparison reflects the combined effect of the random subspace and gradient-balanced weighting.

Under DG-SAM, ADSD achieves the best R@1 and AP on SUES-200 and both GTA-UAV splits among the tested distillation objectives, including TG-HND, full-dimensional alignment, and VkD. On GTA cross-area, ADSD reaches 43.48\% R@1 compared with 39.58\% for VkD. Although VkD performs better on the source University-1652 dataset, ADSD consistently leads on the zero-shot settings. These results support dual-subspace supervision as an effective way to transfer the intermediate Transformer's knowledge to a compact CNN while improving generalization to unseen datasets. Detailed results and analyses of TS/RS dimensions, multiple seeds, and VkD distillation weights are provided in Section F of the supplementary material.


\par\medskip\noindent\begin{minipage}{\columnwidth}
\captionsetup{type=table}
\centering
\caption{Stage-II ablations of student training. Cells report D$\rightarrow$S R@1/AP (\%). Abbreviations follow Table~\ref{tab:ablation_middle}. Full-dimensional alignment uses a residual MLP.}
\label{tab:ablation_student}
\setlength{\tabcolsep}{1.5pt}
\fontsize{7.5}{9}\selectfont
\begin{tabular*}{\linewidth}{@{\extracolsep{\fill}}lcccc@{}}
\toprule
Variant & U1652 & SUES-200 & GTA-S & GTA-C \\
\midrule
\multicolumn{5}{@{}l}{\textit{(a) Subspaces and optimization}} \\
W/o distillation & 88.52/90.32 & 88.69/90.66 & 34.80/43.02 & 36.00/45.06 \\
TS only & 92.24/93.51 & 92.09/93.50 & 39.24/47.67 & 41.98/50.93 \\
ADSD & 92.94/94.06 & 92.98/94.27 & 38.48/46.66 & 41.00/50.11 \\
TS only + DG-SAM & 92.31/93.54 & 92.90/94.21 & 37.82/45.91 & 40.21/49.30 \\
ADSD + DG-SAM & \textbf{93.20}/\textbf{94.28} & \textbf{93.58}/\textbf{94.80} & \textbf{40.37}/\textbf{48.16} & \textbf{43.48}/\textbf{51.99} \\
\midrule
\multicolumn{5}{@{}l}{\textit{(b) Distillation objectives (all with DG-SAM)}} \\
TG-HND & 89.96/91.62 & 89.44/91.45 & 36.51/44.63 & 38.94/48.02 \\
Full-dim. alignment & 92.88/94.02 & 92.43/93.88 & 36.19/44.57 & 38.81/48.00 \\
VkD~\cite{miles2024vkd} & \textbf{93.90}/\textbf{94.91} & 91.44/92.92 & 37.23/45.57 & 39.58/48.39 \\
ADSD & 93.20/94.28 & \textbf{93.58}/\textbf{94.80} & \textbf{40.37}/\textbf{48.16} & \textbf{43.48}/\textbf{51.99} \\
\bottomrule
\end{tabular*}

\end{minipage}\par\medskip
\par\medskip\noindent\begin{minipage}{\columnwidth}
\captionsetup{type=table}
\centering
\caption{Comparison of one-stage and two-stage transfer paths. $T$, $M$, and $S$ denote the foundation teacher, intermediate teacher, and student, respectively. All paths use ADSD and DG-SAM for student training. In the one-stage $M\rightarrow S$ path, $M$ is trained with retrieval supervision only. Cells report D$\rightarrow$S R@1/AP (\%) as in Table~\ref{tab:ablation_middle}.}
\label{tab:transfer-path}
\setlength{\tabcolsep}{1.5pt}
\fontsize{7.5}{9}\selectfont
\begin{tabular*}{\linewidth}{@{\extracolsep{\fill}}lcccc@{}}
\toprule
Transfer path & U1652 & SUES-200 & GTA-S & GTA-C \\
\midrule
\multicolumn{5}{@{}l}{\textit{(a) One-stage distillation}} \\
$T\rightarrow S$ (ADSD) & 92.53/93.68 & 91.74/93.22 & 38.66/46.68 & 40.55/49.79 \\
$T\rightarrow S$ (+TG-HND) & 91.90/93.26 & 91.30/92.95 & 36.28/44.51 & 37.11/46.51 \\
$M\rightarrow S$ (ADSD) & 92.65/93.80 & 92.08/93.55 & 38.11/46.20 & 39.52/48.57 \\
\midrule
\multicolumn{5}{@{}l}{\textit{(b) Two-stage distillation}} \\
$T\rightarrow M\rightarrow S$ (ours) & \textbf{93.20/94.28} & \textbf{93.58/94.80} & \textbf{40.37/48.16} & \textbf{43.48/51.99} \\
\bottomrule
\end{tabular*}

\end{minipage}\par\medskip
\paragraph{Two-stage transfer path.}Table~\ref{tab:transfer-path} compares our two-stage path with one-stage alternatives. Across University-1652, SUES-200, and both GTA-UAV splits, the two-stage path achieves the best R@1 and AP. On GTA cross-area, it raises R@1 from 40.55\% for direct $T\rightarrow S$ transfer with ADSD to 43.48\%, while adding TG-HND to direct transfer does not close the gap. The foundation-guided intermediate teacher also improves upon its retrieval-only counterpart, which reaches 39.52\% on the same split. These results highlight the value of transferring foundation knowledge to a task-adapted intermediate model before supervising the compact student. By addressing the capacity gap first and the architectural discrepancy subsequently, the two-stage path offers a more effective route for preserving generalizable knowledge under the tested settings.
